\chapter{\texttt{pipeline.atmosphere} --- atmospheric production to Earth surface}
\label{ch:pipeline-atmosphere}

\texttt{pipeline.atmosphere} is the production-to-surface leg for
atmospheric neutrinos: given a tabulated production flux for one produced
particle (e.g.\ Honda/HKKM-style tables, binned in energy, production
height, and zenith/off-axis angle), it selects the requested angular bin
and coherently propagates that produced flavour through the atmosphere to
the Earth surface (\cref{ch:atmosphere}). Continuing the resulting surface
state through the Earth to a detector is a separate step, covered by
\cref{ch:pipeline-atmosphere-earth}.

\section{\texttt{pipeline.atmosphere} --- production selection and surface propagation}

\modulehead{pipeline/atmosphere.py}{\pyfunc{select\_production\_flux} and
\pyfunc{propagate\_atmosphere\_to\_surface}: pick one produced
particle/angle from a flux dataset, then propagate it coherently to the
Earth surface.}

\begin{implbox}
\begin{itemize}
  \item \pyfunc{select\_production\_flux(flux\_data, particle, *,
    alpha\_deg=None, theta\_deg=None, angle\_index=None,
    angle\_mode="theta", angle\_tolerance\_deg=None, context=...)}: given a
    \code{\{particle: \{grids, tables, \ldots\}\}} mapping (as produced by
    the project's flux-table loaders), locates the angular bin closest to
    the requested \code{alpha\_deg}/\code{theta\_deg} (or an explicit
    \code{angle\_index}), and returns a single \code{production} dictionary
    -- the energy/height grids, the selected slice of the height-resolved
    production flux \code{phi\_Eh} and the angle-integrated
    \code{phi\_E\_theta}, the resolved \code{theta\_deg}, and the
    \pyfunc{core.common.neutrino.flavour\_index} of \code{particle} -- that
    every function below consumes. \code{angle\_tolerance\_deg} raises if
    no angular bin lies close enough to the requested angle, instead of
    silently returning a distant one.
  \item \pyclass{AtmosphereSurfaceResult} (frozen dataclass): the
    \code{production} dictionary above, the coherent \code{initial\_state}
    and \code{surface\_states} amplitudes, the full evolutor
    \code{S\_atmosphere}, \code{surface\_probabilities}, the built
    \code{trajectories} (\cref{ch:atmosphere}, \texttt{geometry.py}
    section), and -- only when \code{production["phi\_Eh"]} is present --
    the production-flux-weighted \code{surface\_probability\_height\_averaged},
    \code{surface\_flux\_Eh}, and their optional height/energy-integrated
    reductions.
  \item \pyfunc{propagate\_atmosphere\_to\_surface(production, config, *,
    initial\_flavour=None, trajectory\_steps=200, integrate\_energy=False)}:
    builds the coherent initial state (defaulting to
    \code{production["particle"]} itself, i.e.\ no flavour change at
    production), evaluates \pyfunc{atmosphere\_evolutor}
    (\cref{ch:atmosphere}) on the full energy--height grid at the
    production angle, and applies it to obtain \code{surface\_states} and
    \code{surface\_probabilities}. When the production dictionary carries a
    height-resolved flux, it is additionally combined with the
    probabilities via \pyfunc{core.common.flux.flux\_state} and
    \pyfunc{probability\_weighted\_average}/\pyfunc{flux\_integrated\_coordinate}
    (\cref{ch:core-common}) to give the height-averaged probability and
    height-integrated flux directly, so a caller that only needs the
    energy-resolved surface flux does not have to redo that reduction by
    hand.
\end{itemize}
\end{implbox}

\begin{examplebox}[title={Muon-neutrino production at 60\textdegree{} zenith, propagated to the surface}]
\begin{lstlisting}[style=tpython]
import torch

from tpeanuts.config.propagation import PropagationConfig
from tpeanuts.pipeline.atmosphere import (
    propagate_atmosphere_to_surface,
    select_production_flux,
)
from tpeanuts.util.context import RuntimeContext

context = RuntimeContext.resolve("cpu", torch.float64)
oscillation = PropagationConfig.oscillation_parameters_from_preset(
    "_SM_NUFIT52_NO", context=context,
)
config = PropagationConfig(runtime=context, oscillation=oscillation)

# flux_data: {"numu": {...}, "nue": {...}, ...}, as returned by this
# project's Honda/HKKM-style flux-table loader (see medium.atmosphere).
production = select_production_flux(
    flux_data,
    particle="numu",
    theta_deg=60.0,
    angle_tolerance_deg=1.0,
    context=context,
)

surface = propagate_atmosphere_to_surface(
    production,
    config,
    trajectory_steps=200,
)

print(surface.surface_probabilities.shape)          # (E, h, flavour)
print(surface.surface_flux_Eh.shape)                # (E, h, flavour)
print(surface.surface_flux_height_integrated_E.shape)  # (E, flavour)
\end{lstlisting}
\end{examplebox}
